\subsection{Certified windows and the upper bound 241}
\begin{lemma}[Additional internal window exclusions]\label{lem:windows241}
The singleton counts at consecutive internal vertices of a snake cannot form
any of
\[
(2,2,2),\ (2,2,3),\ (2,2,4),\ (2,3,2),\ (3,2,3),
\quad (4,1,c)\quad(4\le c\le6),
\]
or their reversals. Four consecutive internal counts cannot be
$(2,4,2,3)$ or $(3,2,4,2)$.
\end{lemma}
\begin{proof}[Computer-assisted proof]
The first pattern is excluded by the preceding geometric classification.
For the remaining triples, normalize the surrounding five path vertices
to transition sequences $0,1,2,0$ or $0,1,2,3$. For the quadruples,
the surrounding six vertices reduce, allowing reversal, to
$0,1,2,0,3$, $0,1,2,3,0$, or $0,1,2,3,4$.
Translation and coordinate permutation preserve membership, adjacency and
singleton-square counts; reversal also preserves these counts.

Use the union of radius-three Hamming balls about the internal centres.
Assign each selected-vertex variable its membership in the transformed snake,
and each endpoint variable its membership in the two endpoints. The path
degree, endpoint, square and cycle restrictions are necessary conditions,
as in the preceding finite refutations. At each centre the two incident
path directions cannot occur in a singleton square. Each of the remaining
$\binom72=21$ coordinate pairs has its two neighbouring vertices outside the
snake, so its square is a singleton precisely when its opposite corner is
absent. Thus the prescribed missing-corner bounds express the actual
singleton counts. Auxiliary counters receive their values from these
vertex and endpoint assignments; their clauses follow from counter semantics.

Twenty-two finite certificates cover the triples, range $4\le c\le6$,
quadruples, reversals and canonical shapes. Every input clause is checked
against its geometric or counter rule. Reverse unit propagation verifies
each derived clause until the empty clause is obtained. Consequently the
assignment constructed from a snake cannot exist. Both these Boolean
refutations and their semantic interpretation, including arbitrary-index
normalization, are checked by the Lean kernel. An independent reconstruction
from source confirms the complete argument; no solver assertion is an axiom.
\end{proof}

\begin{lemma}[Internal singleton density]\label{lem:density241}
If a snake has $n\ge3$ internal vertices with counts $h_0,\ldots,h_{n-1}$,
then
\[
\sum_{j=0}^{n-1}h_j\ge3n-4.
\]
\end{lemma}
\begin{proof}
Put $q(t)=\min(t,7)$ and $P(a,b,c)=P_0(q(a),q(b),q(c))$, where
Table~\ref{tab:potential241} specifies $P_0$. All counts are positive;
a neighbour of count one has count at least four. Exclude the triples and
quadruples in Lemma~\ref{lem:windows241}. Direct finite verification on
$\{1,\ldots,7\}$ gives, for every permitted triple and quadruple,
\[
0\le P(a,b,c)\le4,\qquad
9\le a+b+c+P(a,b,c),
\]
\[
3+P(a,b,c)\le d+P(b,c,d).
\]
These are finite integer comparisons, also checked in the kernel.
The truncation is faithful: all excluded entries are at most six,
so a value replaced by seven cannot create an excluded pattern. It
preserves the count-one neighbour condition, and only decreases the
right-hand costs. Hence the inequalities hold for the actual counts.

Start with the inequality for the first three counts, then apply the
four-count inequality successively. The potentials telescope, giving
\[
3n\le \sum_{j=0}^{n-1}h_j+P(h_{n-3},h_{n-2},h_{n-1})
\le\sum_{j=0}^{n-1}h_j+4.
\]
\end{proof}

\begin{center}
\small
\begin{longtable}{ccrrrrrrr}
\caption{Potential $P_0(a,b,c)$: columns give $c=1,\ldots,7$.
Unlisted pairs $(a,b)$ have value zero for every $c$.}\label{tab:potential241}\\
\toprule
$a$ & $b$ & 1 & 2 & 3 & 4 & 5 & 6 & 7\\
\midrule\endfirsthead
\toprule $a$ & $b$ & 1 & 2 & 3 & 4 & 5 & 6 & 7\\\midrule\endhead
1 & 4 & 3 & 2 & 1 & 0 & 0 & 0 & 0\\
1 & 5 & 3 & 2 & 1 & 0 & 0 & 0 & 0\\
1 & 6 & 2 & 1 & 0 & 0 & 0 & 0 & 0\\
1 & 7 & 2 & 1 & 0 & 0 & 0 & 0 & 0\\
2 & 2 & 0 & 0 & 0 & 0 & 1 & 0 & 0\\
2 & 3 & 0 & 0 & 2 & 1 & 0 & 0 & 0\\
2 & 4 & 4 & 2 & 2 & 1 & 0 & 0 & 0\\
2 & 5 & 3 & 2 & 1 & 0 & 0 & 0 & 0\\
2 & 6 & 2 & 1 & 0 & 0 & 0 & 0 & 0\\
2 & 7 & 2 & 1 & 0 & 0 & 0 & 0 & 0\\
3 & 2 & 0 & 0 & 0 & 2 & 1 & 0 & 0\\
3 & 3 & 0 & 3 & 2 & 1 & 0 & 0 & 0\\
3 & 4 & 3 & 2 & 1 & 0 & 0 & 0 & 0\\
3 & 5 & 2 & 1 & 0 & 0 & 0 & 0 & 0\\
3 & 6 & 2 & 1 & 0 & 0 & 0 & 0 & 0\\
3 & 7 & 2 & 1 & 0 & 0 & 0 & 0 & 0\\
4 & 2 & 0 & 0 & 2 & 1 & 0 & 0 & 0\\
4 & 3 & 0 & 3 & 2 & 1 & 0 & 0 & 0\\
4 & 4 & 3 & 2 & 1 & 0 & 0 & 0 & 0\\
4 & 5 & 2 & 1 & 0 & 0 & 0 & 0 & 0\\
4 & 6 & 2 & 1 & 0 & 0 & 0 & 0 & 0\\
4 & 7 & 2 & 1 & 0 & 0 & 0 & 0 & 0\\
5 & 1 & 0 & 0 & 0 & 0 & 1 & 0 & 0\\
5 & 2 & 0 & 3 & 2 & 1 & 0 & 0 & 0\\
5 & 3 & 0 & 2 & 1 & 0 & 0 & 0 & 0\\
5 & 4 & 2 & 1 & 0 & 0 & 0 & 0 & 0\\
5 & 5 & 2 & 1 & 0 & 0 & 0 & 0 & 0\\
5 & 6 & 2 & 1 & 0 & 0 & 0 & 0 & 0\\
5 & 7 & 2 & 1 & 0 & 0 & 0 & 0 & 0\\
6 & 2 & 0 & 2 & 1 & 0 & 0 & 0 & 0\\
6 & 3 & 0 & 1 & 0 & 0 & 0 & 0 & 0\\
6 & 4 & 2 & 1 & 0 & 0 & 0 & 0 & 0\\
6 & 5 & 2 & 1 & 0 & 0 & 0 & 0 & 0\\
6 & 6 & 2 & 1 & 0 & 0 & 0 & 0 & 0\\
6 & 7 & 2 & 1 & 0 & 0 & 0 & 0 & 0\\
7 & 1 & 0 & 0 & 0 & 1 & 0 & 0 & 0\\
7 & 2 & 0 & 2 & 1 & 0 & 0 & 0 & 0\\
7 & 3 & 0 & 1 & 0 & 0 & 0 & 0 & 0\\
7 & 4 & 2 & 1 & 0 & 0 & 0 & 0 & 0\\
7 & 5 & 2 & 1 & 0 & 0 & 0 & 0 & 0\\
7 & 6 & 2 & 1 & 0 & 0 & 0 & 0 & 0\\
7 & 7 & 2 & 1 & 0 & 0 & 0 & 0 & 0\\
\bottomrule\end{longtable}\end{center}


\begin{theorem}[Verified upper bound 241]\label{thm:upper241}
Every snake in $\Q_9$ has at most 241 edges. Therefore
\[
\boxed{191\le a(9)\le241.}
\]
\end{theorem}
\begin{proof}
Let $v$ be the number of vertices. For $v\le4$ the conclusion is immediate.
Otherwise there are $n=v-2\ge3$ internal vertices. Lemma~\ref{lem:density241}
and the two endpoint counts, each at least four, imply
\[
T\ge3(v-2)-4+8=3v-2.
\]
The square-incidence identity and nonnegativity of empty-square counts give
$T+35v\le9214$. Thus $38v\le9216<38\cdot243$, so $v\le242$ and the
edge count is at most 241. The published 191-edge witness supplies the
lower bound. The unconditional upper bound and interval theorem have
both passed an independent source reconstruction and transitive axiom
audit in Lean 4.19.0; the only axioms are propositional extensionality,
classical choice and quotient soundness.
\end{proof}
